\documentclass[10pt,a4paper]{article}
\usepackage[margin=1.8cm]{geometry}
\usepackage{amsmath,amssymb}
\usepackage{booktabs}
\usepackage{tabularx}
\newcolumntype{L}{>{\raggedright\arraybackslash}X}
\usepackage{graphicx}
\usepackage{float}
\usepackage[T1]{fontenc}
\usepackage{microtype}
\usepackage[colorlinks,linkcolor=black,citecolor=black,urlcolor=blue]{hyperref}
\usepackage[dvipsnames]{xcolor}
\usepackage{titlesec}
\definecolor{passgreen}{rgb}{0.05,0.50,0.20}
\definecolor{failred}{rgb}{0.72,0.13,0.13}
\newcommand{\PASS}[1]{\textcolor{passgreen}{\textbf{#1}}}
\newcommand{\FAIL}[1]{\textcolor{failred}{\textbf{#1}}}
\newcommand{\PEND}[1]{\textcolor{gray}{#1}}
\titlespacing*{\section}{0pt}{6pt}{2pt}
\newcommand{\code}[1]{\texttt{#1}}
\newcommand{\op}[1]{\texttt{#1}}
\title{\vspace{-1.4cm}\textbf{Death and inhibition as morphogenetic operators:\\
\normalsize removing tissue, stopping growth, and the gates before either enters the loop}}
\author{}
\date{\vspace{-0.9cm}\small \texttt{tyssue\_ops3d.Apoptosis3D} $\cdot$ \texttt{tyssue\_rd\_ops.Grow3D}
$\cdot$ \texttt{make\_sculpt.py} $\cdot$ \texttt{make\_two\_species.py} $\cdot$
\texttt{log/okuda/sc\_*, ts\_*, tsd\_*} $\cdot$ 11 August 2026}
\setlength{\parindent}{0pt}\setlength{\parskip}{0.55em}
\begin{document}\maketitle\vspace{-1.0cm}

\section*{1.\; The mechanism, and the arithmetic that makes it necessary}

Twenty-nine rounds of this campaign, and six of the one running now, have produced spheres, buds and
lobes. Not one finger. The reason is not search: Route~A swept three operators over six ladders and
Route~B made structural edits every round. It is that \textbf{every operator in the vocabulary
deforms the sheet outward}. Growth inflates, division subdivides, the purse-string necks, and the
growth law is
%
\begin{equation}
\dot v_{\rm eq} \;=\; \texttt{rate}\cdot\bigl(\rho + \mathrm{hill}(a)\bigr),
\qquad \mathrm{hill}(a)=\frac{a^{n}}{a_{\rm sw}^{n}+a^{n}} \in [0,1],
\label{eq:growth}
\end{equation}
%
which is \emph{purely activating}. The fastest a cell grows is set by the activator; the slowest is
$\texttt{rate}\cdot\rho$, and $\rho>0$ is required by premise~P1 --- a tissue that adds no material
is not growing. \textbf{So the tissue between the spots always inflates, and a bulge can never
sharpen into a finger.} Six rounds of broad lobes are what Eq.~\eqref{eq:growth} predicts.

A finger requires the flanks to \emph{stop}, and there are exactly two ways to say stop: remove the
cells, or switch their growth off. Those are the two operators here.

\textbf{Death is a Plexus2 family, not a convenience.} The eight elementary families include
\emph{Die} --- removal of biological entities --- and this vesicle had no implementation of it. The
biology is Monier \emph{et al.} (2015): apoptotic force drives epithelial folding in \emph{Drosophila},
and it is an \emph{inward} deformation, which is the one thing nothing here could produce.
Invagination is one of Okuda's three target morphologies.

\textbf{The arithmetic that decides what can be represented.} An apoptotic cell leaves an epithelium
by losing neighbours until it is a triangle, and a triangle collapses to a point as a T2 transition:
$V-2$, $E-3$, $F-1$, so $\chi=V-E+F$ is unchanged and the sheet stays closed at genus~0. Collapsing
anything with $k>3$ sides would merge $k$ vertices into one and leave a vertex of degree $k$, which a
trivalent sheet \emph{cannot represent} --- a rosette. So death is not a single operator: it is
mark, then shed, then extrude, and the shedding belongs to \op{reconnect\_t1\_3d}. That is not a
hidden coupling; it is the mechanism. A cell leaves an epithelium by losing neighbours.

\section*{2.\; Equations and symbols}

\textbf{Death, as three steps.} A marked cell's target volume decays geometrically and it is
extruded when it passes a fraction of the seed-time median $v_{\rm ref}$:
%
\begin{equation}
V^0_f \leftarrow (1-s)\,V^0_f ,\qquad
\text{extrude when } V^0_f < c\,v_{\rm ref},\qquad
\boxed{\;\tau_{\rm clear}=\frac{\ln c}{\ln(1-s)}\ \text{ticks}\;}
\label{eq:clear}
\end{equation}
%
with $s=\texttt{shrink\_rate}$ and $c=\texttt{critical\_frac}$. Eq.~\eqref{eq:clear} is the whole of
the throughput argument and it is checkable in closed form: $12$ ticks at $s=0.15$, $37$ at $0.05$,
$94$ at $0.02$, $189$ at $0.01$, all with $c=0.15$.

\textbf{Volume is conserved into the survivors}, not discarded. When cell $f$ dies its remaining
volume is shared over its live ring neighbours $\mathcal N(f)$:
%
\begin{equation}
c_g \leftarrow c_g + \frac{V^0_f}{|\{g\in\mathcal N(f): V^0_g \ge c\,v_{\rm ref}\}|\cdot V^0_g},
\qquad g \text{ live}.
\label{eq:conserve}
\end{equation}
%
The live-only restriction in the denominator is load-bearing and was found by a failure: dividing by
$\max(V^0_g,10^{-9})$ let a marked-but-undying neighbour sit at the floor and produced an activator
of $-1.04\times10^{10}$.

\textbf{A death RATE, not a death set.} Every state-selected rule re-evaluates each frame, so on a
growing tissue marks accumulate into a wave. With $N$ live cells,
%
\begin{equation}
|\{\text{marked}\}| \le \phi N,\qquad \phi=\texttt{max\_mark\_frac},
\label{eq:cap}
\end{equation}
%
and when the cap binds the \emph{worst} candidates are taken, ranked by the mode's own severity. The
geometric modes (\op{band}, \op{cone}, \op{list}) are \textbf{exempt}: a set chosen once has no flux,
and capping it would delete the experiment.

\textbf{Inhibition supplies the zero Eq.~\eqref{eq:growth} lacks}, multiplicatively:
%
\begin{equation}
\dot v_{\rm eq} \;=\; \texttt{rate}\cdot\bigl(\rho+\mathrm{hill}(a_A)\bigr)\cdot
\bigl(1-\mathrm{hill}_{\rm inh}(\hat a_B)\bigr),
\qquad \hat a_B = a_B/\max_f a_B ,
\label{eq:inhib}
\end{equation}
%
so where the second morphogen saturates the cell does not grow \emph{at all}, baseline included.
$\hat a_B$ is normalised to the inhibitor's own maximum for the same reason $a_{\rm sw}$ is a
fraction: an absolute threshold against a field whose scale the chemistry sets is either always on
or always off, and this project has paid for that twice.

\smallskip
{\small
\begin{tabular}{@{}llp{8.6cm}@{}}\toprule
\multicolumn{3}{@{}l}{\textbf{Death} --- Eqs.~\eqref{eq:clear}--\eqref{eq:cap}}\\
$s$ & \code{shrink\_rate} & fraction of its target volume a marked cell loses PER TICK\\
$c$ & \code{critical\_frac} & extrusion threshold, a fraction OF $v_{\rm ref}$ (not of the cell)\\
$v_{\rm ref}$ & reference volume & the SEED-TIME median cell volume; never recomputed, or every
  threshold in the model moves when cells die\\
$\tau_{\rm clear}$ & clearing time & ticks from sentence to extrusion. \emph{Intensive}: a property
  of the pair $(s,c)$, not of the tissue\\
$\phi$ & \code{max\_mark\_frac} & fraction of LIVE cells under sentence at once. Bounds a
  \emph{flux}; the mode chooses who, this chooses how fast\\
$n_{\rm apop}$ & deaths & cumulative count. \emph{Extensive} --- it scales with the tissue, so it is
  never comparable between runs of different size without dividing by $N$\\
\midrule
\multicolumn{3}{@{}l}{\textbf{Growth and inhibition} --- Eqs.~\eqref{eq:growth}, \eqref{eq:inhib}}\\
$a_A,\ a_B$ & activator of species A, B & \emph{per cell}, in \code{chem} columns $(0,1)$ and $(2,3)$\\
$\rho$ & baseline & the growth every cell gets ungated; tip:body is $(\rho+1):\rho$\\
$a_{\rm sw}$, $n$ & gate threshold, sharpness & a FRACTION of the activator's own maximum\\
\code{inhib\_sw}, \code{inhib\_hill} & inhibition threshold, sharpness & fraction of the
  INHIBITOR's own maximum at which growth is half off\\
\code{chan} & species selector & which \code{chem} pair an operator owns. Default 0, so every
  pre-existing spec is unchanged\\
\bottomrule
\end{tabular}}

\section*{3.\; The Plexus2 container}

No new set: death and inhibition act on sets that already exist. What is new is a \emph{family} the
vesicle did not have, and a second species inside one state buffer.

\begin{center}\small
\begin{tabular}{@{}llll@{}}\toprule
\textbf{set} & \textbf{kind} & \textbf{one element is} & \textbf{state it carries}\\\midrule
\code{vertex} & node & a mesh vertex & \code{pos}\\
\code{cell}   & node & one cell of the epithelium & \code{chem} (\textbf{width 4}: two species),
  \code{cen}, \code{area}\\
\multicolumn{4}{@{}l}{\emph{per-face mesh state, carried on the half-edge table:} $V^0_f$,
  $V_{\rm birth}$, $A^0$, $P^0$, \code{age}, \code{ndiv}, \code{alive}, \code{apop\_flag},
  \code{inhib\_frac}}\\
\bottomrule
\end{tabular}
\end{center}

\begin{center}\footnotesize
\begin{tabular}{@{}lllp{4.6cm}p{3.0cm}@{}}\toprule
\textbf{operator} & \textbf{family} & \textbf{acts on} & \textbf{mechanism} & \textbf{its gate}\\\midrule
\op{apoptosis\_3d} & \textbf{Die} & \code{cell} & Eqs.~\eqref{eq:clear}--\eqref{eq:cap}: mark, shrink,
  extrude as a T2 & D1--D9\\
\op{reconnect\_t1\_3d} & Rewire & \code{vertex} & sheds the marked cell's neighbours --- REQUIRED, or
  a cell shrinks forever & D4\\
\op{grow\_3d} & Lateral & \code{cell} & Eq.~\eqref{eq:inhib}, activation $\times$ inhibition & D11--D12\\
\op{cell\_react} & Lateral & \code{cell} & one instance per species, \code{chan} names its pair & D10\\
\op{cell\_diffuse} & Lateral & \code{cell} & per-column diffusivity vector & D10\\
\op{seed\_cell\_rd} & Seed & \code{cell} & seeds into its own columns & D10\\
\bottomrule
\end{tabular}
\end{center}

\textbf{Schedule, and one ordering that is not cosmetic.} Death runs \emph{between growth and the
mechanics}: growth sets the targets, death overrides them for the sentenced cells, and the
relaxation sees both in one tick, so a cell extruded this frame has its hole closed this frame.
Measured: the same composition with death appended after the recorder instead gave grip $0.031$
against $0.110$, and the weakest parent in the set flipped from a collapse to a clean run.

\smallskip
{\footnotesize
\begin{verbatim}
  seed_mesh_3d -> cell_geometry_3d -> seed_cell_rd(chan 0) -> seed_cell_rd(chan 2)
    -> cell_adjacency -> cell_diffuse(0) -> cell_diffuse(2) -> cell_react(0) -> cell_react(2)
    -> shape_to_chem -> grow_3d -> apoptosis_3d -> shape_energy_3d
    -> interface_line_tension_3d -> reconnect_t1_3d -> divide_3d -> topo_snapshot_3d
\end{verbatim}}
\smallskip

There is no \code{substep\_dt} block: this vesicle integrates at one rate, and \op{shape\_energy\_3d}
carries its own \code{relax\_iters} internally.

\section*{4.\; The gates}

A threshold decided \emph{before} the run, and the measured value. The tiers are the paper's:
\textbf{bookkeeping} asks whether the code does what the operator says, \textbf{closed form} whether
the implementation reproduces physics it was given, and \textbf{measurement} whether the model agrees
with something observed in cells. \textbf{Only the third is a statement about biology, and this note
has none of it} --- see \S6.

\smallskip
{\small
\begin{tabular}{@{}llll@{}}\toprule
& \textbf{gate} & \textbf{threshold, set before} & \textbf{measured}\\\midrule
\multicolumn{4}{@{}l}{\textit{tier 1 --- bookkeeping: does the code do what the operator says?}}\\
D1 & every mode fires on a population whose answer is known, and NOT on a
     uniform field & 15/15 selectors & \PASS{19/19} (4 added for the cap)\\
D2 & a topology change permutes PENDING deltas, not just state & activator
     $\ge 0$ & \PASS{$0$} (was $-1.04{\times}10^{10}$)\\
D3 & a dying cell's volume goes to LIVE neighbours only & no divisor below
     $c\,v_{\rm ref}$ & \PASS{P4, P12 cleared}\\
D4 & $\chi=2$ at every frame with death running & exact & \PASS{$2$, all frames}\\
D5 & $n_{\rm apop}$ equals the cells removed & exact, on a death-only
     composition & \PASS{$400-N_{\rm final}$}\\
D6 & the cap admits at most $\phi N$ & $8$ of $100$ proposed at $\phi N=8$ & \PASS{$8$}\\
D7 & when the cap binds, the WORST are taken & the known-worst index & \PASS{took [20]}\\
D8 & the queue must DRAIN before more are sentenced & stays at $\phi N$ &
     \PASS{$4$ of $4$}\\
D9 & a NAMED population is exempt from the cap & all 180 & \PASS{$180$}\\
D10& two species never write each other's columns & $\{0,1\}$ and $\{2,3\}$ &
     \PASS{exact, on a non-uniform field}\\
D11& \code{inhib\_chan} absent leaves an existing run BIT-IDENTICAL & exact &
     \FAIL{NOT RUN} --- \S6\\
D12& growth $\to 0$ as the inhibitor saturates & $<5\%$ of uninhibited &
     \PASS{$0.978\to0.014$}\\
\midrule
\multicolumn{4}{@{}l}{\textit{tier 2 --- closed form: does the implementation reproduce what it was given?}}\\
D13& clearing time equals $\ln c/\ln(1-s)$ & within 1 tick & \FAIL{NOT MEASURED}
     --- predicted 12/37/94/189\\
D14& a collapse is a T2: $V-2$, $E-3$, $F-1$ & exact, else refuse & \PASS{refuses
     $k>3$; closure re-checked}\\
D15& death with no spatial pattern is EROSION, not sculpture & $\Delta$red\_vol
     $<0.05$ & \PASS{$0.957$ at 3460 deaths}\\
D16& a geometric chisel removes exactly its declared population & exact &
     \PASS{5/5 geometries}\\
D17& the rate cap is not the throughput lever; $s$ is & deaths $\propto$
     $1/\tau_{\rm clear}$ & \PASS{$\phi{\times}50\to{\times}1.7$;
     $s{\times}3\to{\times}4.7$}\\
\midrule
\multicolumn{4}{@{}l}{\textit{tier 3 --- measurement: does it agree with something observed in cells?}}\\
D18& apoptotic rate matches an epithelium's & $1$--$3\%$ of cells per day &
     \PEND{BLOCKED: no \code{units:}}\\
D19& clearance time matches an extrusion in vivo & $30$--$90$ min &
     \PEND{BLOCKED: no \code{units:}}\\
D20& the fold from a dying band matches Monier's amplitude & --- &
     \PEND{BLOCKED: no \code{units:}}\\
\bottomrule
\end{tabular}}
\smallskip

\textbf{Seventeen of twenty gates are tier~1 or tier~2, and three are blocked.} That proportion is
worth reporting rather than hiding: this note establishes that the implementation is self-consistent
and that it reproduces the arithmetic it was given. \textbf{It establishes nothing about biology.}

\section*{5.\; The injection gates: what must hold before the loop may use this}

Verification is not permission. The campaign has twice adopted an operator that was correct and
unusable --- \op{rd\_interface\_tension} ran for nineteen rounds at a default that could not fire,
and \op{apoptosis\_3d} destroyed every parent it touched before the rate cap existed. So the gates
below are about the LOOP, not about the code, and every one is measured on the campaign's own best
parents rather than on a substrate chosen for the operator.

\smallskip
{\small
\begin{tabular}{@{}llll@{}}\toprule
& \textbf{gate} & \textbf{threshold, set before} & \textbf{measured}\\\midrule
I1 & breaks no premise on the campaign's own parents & 0 of 6 broken &
     \PASS{$0$ at $\phi=0.005$}\\
I2 & preserves the parent's morphology metrics & protr, grip within 10\% &
     \PASS{within 3--8\%, 6/6}\\
I3 & is not inert: the acted-ledger counts it & $>0$ frames &
     \PASS{acted on all}\\
I4 & is REACHABLE: vocabulary, schedule slot, emitter & all three &
     \PASS{menu row, \code{SCHEDULE\_ORDER}, \code{\_emit}}\\
I5 & its default FIRES: no inert-by-construction default & mode $\ne$ \code{list} &
     \PASS{\code{competition}}\\
I6 & Route A can sweep it: a base carries it & $\ge1$ ladder &
     \PASS{6 ladders, 23 values}\\
I7 & \textbf{the metric that would show its effect DISCRIMINATES} & the target
     metric separates it from the null & \FAIL{see below}\\
I8 & the inhibitor's null is bit-identical (D11) & exact & \FAIL{NOT RUN}\\
\bottomrule
\end{tabular}}
\smallskip

\textbf{I7 is the one that fails, and it is not a detail.} Across all $24$ death and inhibition runs,
\code{protrusion\_aspect\_max\_final} is $0.000$ in $23$ of them and $0.150$ in the twenty-fourth.
The loop's own standing law L3, established independently over rounds r001--r006, says why:
\emph{aspect\_peak and n\_tubes do not discriminate a finger from a lobe}. So the campaign cannot
currently tell whether death has produced the morphology death was added to produce. It CAN tell
that death deforms --- \code{reduced\_volume} $0.807$ and \code{invagination} $0.343$ on
\op{sc\_crown} against $0.959$ and $0.120$ for its own control --- so the honest injection is
\emph{for the deformation targets, not for the tube target}, until a finger metric exists.

\section*{6.\; What the runs measured, and what came back wrong}

\textbf{Where beats how much, by about thirty to one.} \op{sc\_crown} removed $113$ cells --- a
polar cap, chosen once --- and deformed the body more than \op{tsd\_max} did with $3460$:

\smallskip
{\small\begin{tabular}{@{}lrrrrr@{}}\toprule
run & deaths & cells & protr\_f & invag\_f & red\_vol\\\midrule
\op{sc\_crown} (polar cap)   & \textbf{113} & 16056 & \textbf{1.267} & \textbf{0.343} & \textbf{0.814}\\
\op{sc\_waist} (equator band)&  332 & 14486 & 1.174 & 0.328 & 0.807\\
\op{sc\_waist\_purse}        &  336 & 15534 & 1.249 & 0.243 & 0.862\\
\op{tsd\_max} (max rate)     & \textbf{3460} &  2895 & 1.140 & 0.111 & 0.957\\
control (no death)           &    0 & 12692 & 1.095 & 0.120 & 0.959\\
\bottomrule\end{tabular}}
\smallskip

Scattered death erodes a ball into a smaller ball; a chisel carves. This is D15 and D16 as a
positive result rather than a check.

\textbf{The inhibitor works and refutes its own hypothesis.} Monotone in strength --- $12692 \to
6920 \to 4859 \to 4154$ cells as \code{inhib\_sw} tightens --- with \code{protr\_final} falling
$1.095 \to 1.020$. I argued in \S1 that a finger needs a zero at the flanks. The zero is there, and
there is no finger: inhibition makes a smaller ball, not a sharper one. The missing ingredient is
therefore \emph{not} the absence of a stop, and \S1's argument is wrong as an explanation even though
it is right as arithmetic.

\textbf{The two-species claim was killed by its own control.} \op{ts\_death\_b\_late} (death on the
second map) against \op{ts\_death\_same\_a} (death on the growth map): $301$ vs $195$ deaths,
\code{protr\_final} $1.082$ vs $\mathbf{1.106}$. The same-species control is marginally BETTER, so
the second map bought nothing. That control existed to be able to say this.

\textbf{Three failures worth keeping.} \op{tsd\_cap10} and \op{tsd\_cap25} broke P4 and P12
(chemistry extinguished, then NaN), so their numbers are not evidence. \op{sc\_old\_core} --- kill the
oldest fastest --- broke P1 and crushed the tissue to $1585$ cells with negative grip: removing the
body outruns growth. And \textbf{the slow-death series is confounded by my own construction}: every
variant carries \code{inhib\_sw} $0.35$, which alone takes the tissue from $12692$ to $4859$ cells, so
those five runs measure inhibition and not death speed. The question they were built to answer is
still open.

\section*{7.\; What this licenses, and what it does not}

\textbf{Licensed.} The Die family exists and is verified at the level of bookkeeping: selectors fire
where they should and nowhere else (D1), a topology change carries its pending deltas (D2), volume is
conserved into survivors (D3), the sheet stays closed (D4), the count is exact (D5), and the rate is
a rate (D6--D9). Two species coexist without leaking (D10). Growth can be switched off (D12). Death
patterned in SPACE deforms the body at a thirtieth of the cell cost of death patterned by rate
(D15--D17), which is the result worth carrying forward.

\textbf{Not licensed.} No gate here is a measurement (D18--D20 blocked for want of a \code{units:}
block, which is the same defect the \code{ecm} study found turned ``$0.82$ grid cells'' into
$15\,\mu$m). The clearing time is predicted and not verified (D13). The inhibitor's null has not
been run (D11/I8), so ``off by default changes nothing'' is an assertion. And I7 fails: the campaign
cannot yet distinguish a finger from a lobe, so it cannot be told to search for one. Injecting for
the \emph{deformation} seats --- where \code{reduced\_volume} and \code{invagination} do
discriminate, by a factor of two against the control --- is supported today; injecting for the tube
seat is not.

\section*{8.\; References}

{\small
\begin{thebibliography}{9}\setlength{\itemsep}{0pt}
\bibitem{monier2015} Monier, B., Gettings, M., Gay, G., Mangeat, T., Schott, S., Guarner, A.,
Suzanne, M. (2015) Apico-basal forces exerted by apoptotic cells drive epithelium folding.
\emph{Nature} \textbf{518}:245--248. --- the mechanism of \S1, and the only inward one available.
\bibitem{okuda2013} Okuda, S., Inoue, Y., Eiraku, M., Sasai, Y., Adachi, T. (2013) Reversible
network reconnection model for simulating large deformation in dynamic tissue morphogenesis.
\emph{Biomech. Model. Mechanobiol.} \textbf{12}(4):627--644. --- the T1/T2 machinery of \S1.
\bibitem{okuda2018} Okuda, S., Miura, T., Inoue, Y., Adachi, T., Eiraku, M. (2018) Combining
Turing and 3D vertex models reproduces autonomous multicellular morphogenesis.
\emph{Sci. Rep.} \textbf{8}:2386. --- the three target morphologies.
\bibitem{gray1984} Gray, P., Scott, S. K. (1984) Autocatalytic reactions in the isothermal,
continuous stirred tank reactor. \emph{Chem. Eng. Sci.} \textbf{39}:1087--1097. --- both species.
\bibitem{gm1972} Gierer, A., Meinhardt, H. (1972) A theory of biological pattern formation.
\emph{Kybernetik} \textbf{12}:30--39. --- lateral inhibition, which Eq.~\eqref{eq:inhib} is a
crude form of.
\end{thebibliography}}

\end{document}
